% --- LaTeX Homework Template - S. Venkatraman ---

% --- Set document class and font size ---

\documentclass[letterpaper, 11pt]{article}

% --- Package imports ---

\usepackage{
  amsmath, amsthm, amssymb, mathtools, dsfont,	  % Math typesetting
  graphicx, wrapfig, subfig, float,                  % Figures and graphics formatting
  listings, color, inconsolata, pythonhighlight,     % Code formatting
  fancyhdr, sectsty, hyperref, enumerate, enumitem } % Headers/footers, section fonts, links, lists

% --- Page layout settings ---

% Set page margins
\usepackage[left=0.7in, right=0.7in, bottom=1in, top=1.1in, headsep=0.2in]{geometry}

% Anchor footnotes to the bottom of the page
\usepackage[bottom]{footmisc}



% Set line spacing
\renewcommand{\baselinestretch}{1.2}

% Set spacing between paragraphs
\setlength{\parskip}{1.5mm}

% Allow multi-line equations to break onto the next page
\allowdisplaybreaks

% Enumerated lists: make numbers flush left, with parentheses around them
\setlist[enumerate]{wide=0pt, leftmargin=21pt, labelwidth=0pt, align=left}
\setenumerate[1]{label={(\arabic*)}}

% --- Page formatting settings ---

% Set link colors for labeled items (blue) and citations (red)
\hypersetup{colorlinks=true, linkcolor=blue, citecolor=red}

% Make reference section title font smaller
\renewcommand{\refname}{\large\bf{References}}

% --- Settings for printing computer code ---

% Define colors for green text (comments), grey text (line numbers),
% and green frame around code
\definecolor{greenText}{rgb}{0.5, 0.7, 0.5}
\definecolor{greyText}{rgb}{0.5, 0.5, 0.5}
\definecolor{codeFrame}{rgb}{0.5, 0.7, 0.5}

% Define code settings
\lstdefinestyle{code} {
  frame=single, rulecolor=\color{codeFrame},            % Include a green frame around the code
  numbers=left,                                         % Include line numbers
  numbersep=8pt,                                        % Add space between line numbers and frame
  numberstyle=\tiny\color{greyText},                    % Line number font size (tiny) and color (grey)
  commentstyle=\color{greenText},                       % Put comments in green text
  basicstyle=\linespread{1.1}\ttfamily\footnotesize,    % Set code line spacing
  keywordstyle=\ttfamily\footnotesize,                  % No special formatting for keywords
  showstringspaces=false,                               % No marks for spaces
  xleftmargin=1.95em,                                   % Align code frame with main text
  framexleftmargin=1.6em,                               % Extend frame left margin to include line numbers
  breaklines=true,                                      % Wrap long lines of code
  postbreak=\mbox{\textcolor{greenText}{$\hookrightarrow$}\space} % Mark wrapped lines with an arrow
}

% Set all code listings to be styled with the above settings
\lstset{style=code}

% --- Math/Statistics commands ---

% Add a reference number to a single line of a multi-line equation
% Usage: "\numberthis\label{labelNameHere}" in an align or gather environment
\newcommand\numberthis{\addtocounter{equation}{1}\tag{\theequation}}

% Shortcut for bold text in math mode, e.g. $\b{X}$
\let\b\mathbf

% Shortcut for bold Greek letters, e.g. $\bg{\beta}$
\let\bg\boldsymbol

% Shortcut for calligraphic script, e.g. %\mc{M}$
\let\mc\mathcal

% \mathscr{(letter here)} is sometimes used to denote vector spaces
\usepackage[mathscr]{euscript}

% Convergence: right arrow with optional text on top
% E.g. $\converge[w]$ for weak convergence
\newcommand{\converge}[1][]{\xrightarrow{#1}}

% Normal distribution: arguments are the mean and variance
% E.g. $\normal{\mu}{\sigma}$
\newcommand{\normal}[2]{\mathcal{N}\left(#1,#2\right)}

% Uniform distribution: arguments are the left and right endpoints
% E.g. $\unif{0}{1}$
\newcommand{\unif}[2]{\text{Uniform}(#1,#2)}

% Independent and identically distributed random variables
% E.g. $ X_1,...,X_n \iid \normal{0}{1}$
\newcommand{\iid}{\stackrel{\smash{\text{iid}}}{\sim}}

% Equality: equals sign with optional text on top
% E.g. $X \equals[d] Y$ for equality in distribution
\newcommand{\equals}[1][]{\stackrel{\smash{#1}}{=}}

% Math mode symbols for common sets and spaces. Example usage: $\R$
\newcommand{\R}{\mathbb{R}}   % Real numbers
\newcommand{\C}{\mathbb{C}}   % Complex numbers
\newcommand{\Q}{\mathbb{Q}}   % Rational numbers
\newcommand{\Z}{\mathbb{Z}}   % Integers
\newcommand{\N}{\mathbb{N}}   % Natural numbers
\newcommand{\F}{\mathcal{F}}  % Calligraphic F for a sigma algebra
\newcommand{\El}{L} % Calligraphic L, e.g. for L^p spaces

% Math mode symbols for probability
\newcommand{\pr}{\mathbb{P}}    % Probability measure
\newcommand{\E}{\mathbb{E}}     % Expectation, e.g. $\E(X)$
\newcommand{\var}{\text{Var}}   % Variance, e.g. $\var(X)$
\newcommand{\cov}{\text{Cov}}   % Covariance, e.g. $\cov(X,Y)$
\newcommand{\corr}{\text{Corr}} % Correlation, e.g. $\corr(X,Y)$
\newcommand{\B}{\mathcal{B}}    % Borel sigma-algebra

% Other miscellaneous symbols
\newcommand{\tth}{\text{th}}	% Non-italicized 'th', e.g. $n^\tth$
\newcommand{\Oh}{\mathcal{O}}	% Big-O notation, e.g. $\O(n)$
\newcommand{\1}{\mathds{1}}	% Indicator function, e.g. $\1_A$

% Additional commands for math mode
\DeclareMathOperator*{\argmax}{argmax}    % Argmax, e.g. $\argmax_{x\in[0,1]} f(x)$
\DeclareMathOperator*{\argmin}{argmin}    % Argmin, e.g. $\argmin_{x\in[0,1]} f(x)$
\DeclareMathOperator*{\spann}{Span}       % Span, e.g. $\spann\{X_1,...,X_n\}$
\DeclareMathOperator*{\bias}{Bias}        % Bias, e.g. $\bias(\hat\theta)$
\DeclareMathOperator*{\ran}{ran}          % Range of an operator, e.g. $\ran(T) 
\DeclareMathOperator*{\dv}{d\!}           % Non-italicized 'with respect to', e.g. $\int f(x) \dv x$
\DeclareMathOperator*{\diag}{diag}        % Diagonal of a matrix, e.g. $\diag(M)$
\DeclareMathOperator*{\trace}{trace}      % Trace of a matrix, e.g. $\trace(M)$

% Numbered theorem, lemma, etc. settings - e.g., a definition, lemma, and theorem appearing in that 
% order in Section 2 will be numbered Definition 2.1, Lemma 2.2, Theorem 2.3. 
% Example usage: \begin{theorem}[Name of theorem] Theorem statement \end{theorem}
\theoremstyle{definition}
\newtheorem{theorem}{Theorem}[section]
\newtheorem{proposition}[theorem]{Proposition}
\newtheorem{lemma}[theorem]{Lemma}
\newtheorem{corollary}[theorem]{Corollary}
\newtheorem{definition}[theorem]{Definition}
\newtheorem{example}[theorem]{Example}
\newtheorem{remark}[theorem]{Remark}

% Un-numbered theorem, lemma, etc. settings
% Example usage: \begin{lemma*}[Name of lemma] Lemma statement \end{lemma*}
\newtheorem*{theorem*}{Theorem}
\newtheorem*{proposition*}{Proposition}
\newtheorem*{lemma*}{Lemma}
\newtheorem*{corollary*}{Corollary}
\newtheorem*{definition*}{Definition}
\newtheorem*{example*}{Example}
\newtheorem*{remark*}{Remark}
\newtheorem*{claim}{Claim}
\newcommand{\Tau}{\mathcal{T}}
\newcommand{\PSet}{\mathcal{P}}


% --- Left/right header text (to appear on every page) ---

% Include a line underneath the header, no footer line
\pagestyle{fancy}
\renewcommand{\footrulewidth}{0pt}
\renewcommand{\headrulewidth}{0.4pt}

% Left header text: course name/assignment number
\lhead{MAM4000W, Analysis II -- Assignment 3}

% Right header text: your name
\rhead{Alex Cristaudo, CRSALE010}

% --- Document starts here ---

\begin{document}
\subsection*{Question 1}
\begin{enumerate}
	\item[(a)] Let $(x_n) \in V$. Then there exists $N \in \mathbb{N}$ such that $x_n = 0$ for all $n \ge N$. Then \[
	\sup \limits _{n \in \mathbb{N}}|T_n(x)| = \sup \limits _{n \in \mathbb{N}} |n| \cdot || e^n ||\cdot  |x_n|
	\]
	For $n \ge N$, $|x_n| = 0$, so it follows that \[
	\sup \limits _{n \in \mathbb{N}}|T_n(x)| = \sup \limits _{\substack{n \in \mathbb{N} \\ n \le N}} |n| \cdot ||e^n || \cdot |x_n| \le N \max \limits _{\substack{n \in \mathbb{N} \\ n \le N}} (||e^n||) \max \limits _{\substack{n \in \mathbb{N} \\ n \le N}} |x_n|
	\]
	Since $||e^n||$ is finite, as $||\cdot ||$ is a norm on $V$, these are maximums of finite sets, so they exist and are finite. Thus \[
	\sup \limits _{n\in \mathbb{N}} |T_n (x)| < \infty
	\]
	\item[(b)] Suppose, for a contradiction, $(V, ||\cdot ||)$ is Banach. We have that each $T_n \in V'$ so each $T_n$ is a linear continuous mapping. Thus, by the Banach-Steinhaus theorem, \[
	\sup \limits _{n \in \mathbb{N}} ||T_n||_{V'} < \infty
	\]
  For any $x = (x_n) \in V$, we have that \[
  |T_n(x)| \le ||x||_V \cdot ||T_n ||_{V'}
  \]
So, for $x \ne 0$,
  \[
  \sup \limits _{n \in \mathbb{N}} ||T_n||_{V'}
 \ge \sup \limits _{n \in \mathbb{N}} \frac{|T_n(x)|}{||x||} 
 \]
 This is true for every $x \in V \setminus \{0\}$. But now, for each $N \in \mathbb{N}$, for $x = (e^N)$, \[
\frac{|T_N((e^N))|}{||(e^N)||} = \frac{|N| \cdot ||e^N|| \cdot |e_N^N|}{||e^N||} = N
\]
But this implies \[
\sup \limits _{n \in \mathbb{N}} ||T_n||_{V'} = \infty
\]
a contradiction. So $(V, ||\cdot ||)$ is not Banach.
\end{enumerate}

\subsection*{Question 2}
\begin{enumerate}
\item[(a)] First notice, for $\alpha, \beta \in \mathbb{K}$ and $L_1, L_2 \in V'$; \[
\Phi(\alpha L_1 + \beta L_2) = (\alpha L_1 + \beta L_2)|_M = \alpha L_1 |_M + \beta L_2|_M = \alpha \Phi(L_1) + \beta \Phi (L_2)
\]
So $\Phi$ is linear. On the other hand, \[
||\Phi(L)||_{M'} = || L|_M ||_{M'} = \sup \limits _{\substack{||v||_M = 1 \\ v \in M}} ||L(v)||_M = \sup \limits _{\substack{||v||_M = 1 \\ v \in M}} ||L(v)||_V \le  \sup \limits _{\substack{||v||_V = 1 \\ v \in V}} ||L(v)||_V  = ||L||_{V'}
\]
So $\Phi$ is continuous and $||\Phi||_{\mathcal{L}(V', M')} \le 1$. Now let $\ell: M \to \mathbb{K}$ be a nonzero linear continuous map. Then, by the Hahn Banach Theorem, there is a linear continuous map $L: V \to \mathbb{K}$ such that \[
L(v) = \ell (v), \ \ \ \forall v \in M, \quad \text{and} \quad ||L||_{V'} = ||\ell||_{M'}
\]
It thus follows that $\Phi (L) = \ell$, so \[
||\Phi||_{\mathcal{L}(V', M')} \ge \frac{||\Phi(L)||_{M'}}{||L||_{V'}} = \frac{||\ell ||_M'}{||L||_{V'}} = 1
\]
Thus $||\Phi ||_{\mathcal{L}(V', M')} = 1$

\item[(b)] Suppose $\Phi$ is injective. Assume that $\overline{M} \ne V$. Then $\overline{M}$ is a proper closed subspace of $V$. Thus for any $v \in V \setminus \overline{M}$, there exists some $L_v \in V'$ such that \[
L_v(u) = 0, \ \ \forall u \in \overline{M}, \quad L_v(v) > 0 
\]
But then $0_{V'} \in V'$ and \[
\Phi(0_{V'}) = 0_{M'}, \quad \Phi(L_v) = 0_{M'}, \quad L_v \ne 0_{V'}
\]
which contradicts that $\Phi$ is injective. So $\overline{M} = V$ and $M$ is dense in $V$. ~\\ ~\\On the other hand, suppose $M$ is dense in $V$ and that $\Phi(L_1) = \Phi(L_2)$. Let $v \in V$. Then there is a sequence $(m_n)$ in $M$ converging to $v$. Thus \[
L_1(v) = L_1 \left (\lim \limits _{n \to \infty} m_n \right) = \lim \limits _{n \to \infty} L_1(m_n)
\]
Using the fact that $L_1$ is continuous. Then $\Phi(L_1) = \Phi(L_2)$ means $L_1(m_n) = L_2(m_n)$ and so \[
\lim \limits _{n \to \infty} L_1(m_n) = \lim \limits _{n \to \infty} L_2(m_n) = L_2 \left (\lim \limits _{n \to \infty} m_n \right) = L_2 (v)
\]
Thus $L_1 = L_2$ and $\Phi$ is injective.

\item [(c)] Let $\ell \in M'$. By the Hahn Banach Theorem, there exists $L \in V'$ such that $L(v) = \ell (v)$ for all $v \in M$. But then $L|_M = \ell$ and so $\Phi(L) = \ell$. Thus $\Phi$ is surjective. Now $V'$ and $M'$ are both Banach spaces over the same field (dual spaces are always Banach), so, since we have shown $\Phi$ is surjective, by the Open Mapping Theorem, $\Phi$ is open.
\end{enumerate}

\end{document}
